\subsubsection{解法の解説}
  満を持して特性方程式の解が虚数解を有する場合を考える.
  ここで,これまで扱ってきた隣接3\,項間漸化式
  \[
    a_{n+2}+pa_{n+1}+qa_n=0
  \]
  を考える.特性方程式は
  \[
    x^2+px+q=0
  \]
  である.

  これまで,特性方程式の解が異なる実数
  \[
    x=\alpha,\beta
  \]
  をもつ場合には,
  \[
    a_n=A\alpha^n+B\beta^n
  \]
  という形の一般項が得られることを学んだ(3.5.1では指数が$n-1$であったが,\,$A,B$は適当な定数のため自由に$\alpha,\beta$を掛けても良い).

  では,特性方程式が実数解をもたず,
  \[
    x=\alpha,\overline{\alpha}
  \]
  という共役な複素数解をもつ場合には,一般項はどのような形になるだろうか.

  例えば,
  \[
    x^2-2x+2=0
  \]
  を考えると,
  \[
    x=1\pm i
  \]
  である.したがって,漸化式
  \[
    a_{n+2}=2a_{n+1}-2a_n
  \]
  の特性方程式は虚数解をもつ.

  形式的には,一般項を
  \[
    a_n=A(1+i)^n+B(1-i)^n
  \]
  と書くことができる.しかし,\,$a_n$が実数列である場合,このままでは一般項の実数性が分かりにくい.そこで,複素数の極形式を利用する.

  一般に,特性方程式の解を
  \[
    \alpha=\rho(\cos\theta+i\sin\theta)=\rho e^{i\theta}
  \]
  とすると,共役解は
  \[
    \overline{\alpha}=\rho(\cos\theta-i\sin\theta)=\rho e^{-i\theta}
  \]
  である.したがって,
  \[
    \alpha^n=\rho^n e^{in\theta}
    =\rho^n(\cos n\theta+i\sin n\theta)
  \]
  および
  \[
    \overline{\alpha}^{\,n}
    =\rho^n e^{-in\theta}
    =\rho^n(\cos n\theta-i\sin n\theta)
  \]
  となる.

  この2つを組み合わせると,
  \[
    A\alpha^n+B\overline{\alpha}^{\,n}
  \]
  は,適切な実数定数\,$C,D$を用いて
  \[
    a_n=\rho^n(C\cos n\theta+D\sin n\theta)
  \]
  という実数の形に書き直すことができる.

  したがって,虚数解をもつ隣接3\,項間漸化式では,
  \[
    a_n=\rho^n(C\cos n\theta+D\sin n\theta)
  \]
  という形が基本的な一般項となる.ここで,\,$\rho$は特性根の絶対値,\,$\theta$はその偏角である.

  では,一般項
  \[
    a_n=\rho^n(C\cos n\theta+D\sin n\theta)
  \]
  に含まれる定数$C,D$をどのように決定すればよいだろうか.

  これまでの漸化式と同様に,初期条件を代入することで$C,D$を決定することができる.例えば,\,$a_1,a_2$が与えられているとする.一般項に$n=1,2$を代入すると,
  \[
    a_1=\rho(C\cos\theta+D\sin\theta)
  \]
  および
  \[
    a_2=\rho^2(C\cos2\theta+D\sin2\theta)
  \]
  を得る.

  これは$C,D$についての連立方程式である.したがって,この2つの方程式を解けば$C,D$が決定される.

  特に,\,$a_0$が初期条件として与えられている場合には,より簡単になる.一般項に$n=0$を代入すると,
  \[
    a_0=C
  \]
  となる.したがって,
  \[
    C=a_0
  \]
  が直ちに得られる.

  さらに,\,$n=1$を代入すると,
  \[
    a_1=\rho(C\cos\theta+D\sin\theta)
  \]
  であるから,
  \[
    D=
    \frac{a_1/\rho-C\cos\theta}{\sin\theta}
  \]
  となる.ここで,\,$\sin\theta\neq0$であることに注意する.特性根が実数ではなく虚部をもつ場合には,\,$\theta$は$0$や$\pi$ではないため,
  \[
    \sin\theta\neq0
  \]
  が成り立つ.

  したがって,\,$a_0,a_1$が与えられている場合には,
  \[
    C=a_0,\qquad
    D=\frac{a_1/\rho-a_0\cos\theta}{\sin\theta}
  \]
  として$C,D$を求めることができる.

  例えば,先ほどの
  \[
    a_{n+2}=2a_{n+1}-2a_n
  \]
  を考える.特性根は
  \[
    1+i
    =\sqrt{2}
    \left(
      \cos\frac{\pi}{4}
      +i\sin\frac{\pi}{4}
    \right)
  \]
  であるから,
  \[
    \rho=\sqrt{2},\qquad
    \theta=\frac{\pi}{4}
  \]
  である.したがって一般項は
  \[
    a_n=(\sqrt{2})^n
    \left(
      C\cos\frac{n\pi}{4}
      +D\sin\frac{n\pi}{4}
    \right)
  \]
  となる.

  例えば初期条件
  \[
    a_0=1,\qquad a_1=2
  \]
  が与えられているとする.まず,
  \[
    C=a_0=1
  \]
  である.また,
  \[
    2=\sqrt{2}
    \left(
      C\cos\frac{\pi}{4}
      +D\sin\frac{\pi}{4}
    \right)
  \]
  であるから,
  \[
    2=\sqrt{2}
    \left(
      \frac{1}{\sqrt{2}}
      +\frac{D}{\sqrt{2}}
    \right)
    =1+D
  \]
  より,
  \[
    D=1
  \]
  を得る.

  よって,この場合の一般項は
  \[
    a_n=(\sqrt{2})^n
    \left(
      \cos\frac{n\pi}{4}
      +\sin\frac{n\pi}{4}
    \right)
  \]
  となる.

  このように,虚数解をもつ場合でも,基本的な流れは実数解の場合と変わらない.まず特性根から$\rho,\theta$を求め,
  \[
    a_n=\rho^n(C\cos n\theta+D\sin n\theta)
  \]
  とおき,初期条件を代入して$C,D$を決定すればよい.

  ここまででは,特性方程式の解を利用して一般項の形を導いた.実は,特性方程式の解が複素数である場合には,漸化式そのものを変形することによっても,
  \[
    a_n=\rho^n(C\cos n\theta+D\sin n\theta)
  \]
  という形を自然に導くことができる.

  特性方程式の解を極形式を用いて
  \[
    \alpha=\rho(\cos\theta+i\sin\theta)
  \]
  とすると,共役な解は
  \[
    \overline{\alpha}
    =\rho(\cos\theta-i\sin\theta)
  \]
  である.解と係数の関係より,
  \[
    \alpha+\overline{\alpha}=-p
  \]
  および
  \[
    \alpha\overline{\alpha}=q
  \]
  であるから,
  \[
    p=-2\rho\cos\theta,\qquad
    q=\rho^2
  \]
  を得る.

  したがって,もとの漸化式
  \[
    a_{n+2}+pa_{n+1}+qa_n=0
  \]
  は,
  \[
    a_{n+2}
    -2\rho\cos\theta\,a_{n+1}
    +\rho^2a_n=0
  \]
  と書き直すことができる.

  ここで,両辺を$\rho^{n+2}$で割ると,
  \[
    \frac{a_{n+2}}{\rho^{n+2}}
    -2\cos\theta
    \frac{a_{n+1}}{\rho^{n+1}}
    +\frac{a_n}{\rho^n}
    =0
  \]
  となる.

  そこで,
  \[
    b_n=\frac{a_n}{\rho^n}
  \]
  とおく.すると,
  \[
    b_{n+2}
    -2\cos\theta\,b_{n+1}
    +b_n=0
  \]
  すなわち,
  \[
    b_{n+2}+b_n
    =2\cos\theta\,b_{n+1}
  \]
  を得る.

  この式は,三角関数の加法定理
  \[
    \cos(x+\theta)+\cos(x-\theta)
    =2\cos\theta\cos x
  \]
  および
  \[
    \sin(x+\theta)+\sin(x-\theta)
    =2\cos\theta\sin x
  \]
  と同じ形をしている.

  実際,
  \[
    b_n=\cos n\theta
  \]
  とすると,
  \[
  \begin{aligned}
    b_{n+2}+b_n
    &=
    \cos(n+2)\theta+\cos n\theta
  \\
    &=
    2\cos\theta\cos(n+1)\theta
  \\
    &=
    2\cos\theta\,b_{n+1}
  \end{aligned}
  \]
  となり,確かに漸化式を満たす.

  同様に,
  \[
    b_n=\sin n\theta
  \]
  とすると,
  \[
  \begin{aligned}
    b_{n+2}+b_n
    &=
    \sin(n+2)\theta+\sin n\theta
  \\
    &=
    2\cos\theta\sin(n+1)\theta
  \\
    &=
    2\cos\theta\,b_{n+1}
  \end{aligned}
  \]
  となる.

  したがって,この漸化式では
  \[
    \cos n\theta
  \]
  と
  \[
    \sin n\theta
  \]
  が基本的な解として現れる.

  漸化式は2階の漸化式であるから,初期条件を2つ与えることで数列が一意に定まる.そこで,この2つの基本的な解を組み合わせると,
  \[
    b_n=C\cos n\theta+D\sin n\theta
  \]
  と表すことができる.

  もともと
  \[
    b_n=\frac{a_n}{\rho^n}
  \]
  とおいていたので,
  \[
    a_n=\rho^n b_n
  \]
  である.したがって,
  \[
    a_n=\rho^n
    (C\cos n\theta+D\sin n\theta)
  \]
  を得る.

  このように,複素数解を直接べき乗することなく,特性根から得られる
  \[
    p=-2\rho\cos\theta,\qquad q=\rho^2
  \]
  を用いて漸化式を変形すると,三角関数の加法定理が現れ,そこから
  \[
    \rho^n\cos n\theta,\qquad
    \rho^n\sin n\theta
  \]
  という2つの基本的な形が自然に導かれる.

  ここまでの議論から分かるように,虚数解をもつ漸化式の一般項には
  \[
    \cos n\theta,\qquad \sin n\theta
  \]
  という三角関数が現れる.このことから,数列$\{a_n\}$は三角関数に由来する周期的な振る舞いをもつことが分かる.

  特に,
  \[
    \rho=1
  \]
  の場合には,
  \[
    a_n=C\cos n\theta+D\sin n\theta
  \]
  となるため,数列そのものが周期的になる場合がある.

  実際,
  \[
    \frac{\theta}{2\pi}
  \]
  が有理数であるとする.このとき,ある正の整数$T$と整数$k$を用いて
  \[
    T\theta=2k\pi
  \]
  と書くことができる.したがって,
  \[
    \cos(n+T)\theta=\cos n\theta
  \]
  および
  \[
    \sin(n+T)\theta=\sin n\theta
  \]
  であるから,
  \[
    a_{n+T}=a_n
  \]
  となる.よって,この場合には数列$\{a_n\}$は周期的である.

  一方,\,$\rho\neq1$の場合には,
  \[
    a_n=\rho^n
    (C\cos n\theta+D\sin n\theta)
  \]
  であるため,三角関数の部分は周期的に振る舞うものの,\,$\rho^n$によって数列全体の大きさが変化する.したがって,一般には数列そのものは周期的ではない.しかし,その振動には$\cos n\theta$や$\sin n\theta$に由来する周期的な性質が現れている.

  さらに,このような三角関数の性質に着目すると,一般項をあらかじめ仮定することなく,数学的帰納法によってその形を確かめることもできる.

  例えば,
  \[
    b_{n+2}-2\cos\theta\,b_{n+1}+b_n=0
  \]
  を考える.三角関数の加法定理より,
  \[
    \cos(n+2)\theta+\cos n\theta
    =2\cos\theta\cos(n+1)\theta
  \]
  であるから,
  \[
    \cos(n+2)\theta
    =2\cos\theta\cos(n+1)\theta-\cos n\theta
  \]
  が成り立つ.

  ここで,
  \[
    b_0=1,\qquad b_1=\cos\theta
  \]
  とすると,
  \[
    b_n=\cos n\theta
  \]
  が成り立つことを数学的帰納法によって示すことができる.

  実際,\,$b_n=\cos n\theta$および$b_{n+1}=\cos(n+1)\theta$が成り立つと仮定すると,漸化式より,
  \[
  \begin{aligned}
    b_{n+2}
    &=2\cos\theta\,b_{n+1}-b_n\\
    &=2\cos\theta\cos(n+1)\theta-\cos n\theta\\
    &=\cos(n+2)\theta
  \end{aligned}
  \]
  となる.したがって,数学的帰納法により,
  \[
    b_n=\cos n\theta
  \]
  がすべての$n$について成り立つ.

  同様に,
  \[
    \sin(n+2)\theta
    =2\cos\theta\sin(n+1)\theta-\sin n\theta
  \]
  であるから,
  \[
    b_0=0,\qquad b_1=\sin\theta
  \]
  に対して,
  \[
    b_n=\sin n\theta
  \]
  が数学的帰納法によって示される.

  このように,三角関数の加法定理を利用することで,
  \[
    \cos n\theta,\qquad \sin n\theta
  \]
  が漸化式を満たすことを直接確認できる.さらに,初期条件によってこれらを適切に組み合わせることで,
  \[
    a_n=\rho^n
    (C\cos n\theta+D\sin n\theta)
  \]
  という形が得られる.

  したがって,虚数解をもつ漸化式では,特性方程式から形式的に一般項を求めるだけでなく,三角関数の周期性や加法定理に着目することで,数列がもつ振動的な性質を直接捉えることができる.

  これらの議論を踏まえると,特性方程式の解が実数解の場合には$\alpha^n$や$\beta^n$による指数的な増減が中心となるのに対し,虚数解の場合には
  \[
    \cos n\theta,\qquad\sin n\theta
  \]
  が現れるため,数列に周期的な振動が生じる.さらに,\,$\rho$の大きさによってその振幅が変化する.このとき,\,$\rho$によって
  \[
    \begin{cases}
      \rho<1 & \text{振幅が減衰する}\\
      \rho=1 & \text{振幅が一定である}\\
      \rho>1 & \text{振幅が増大する}
    \end{cases}
  \]
  という違いが生じる.

  とくに,
  \[
    \rho=1
  \]
  の場合には,
  \[
    a_n=C\cos n\theta+D\sin n\theta
  \]
  となり,数列の振る舞いは$\theta$によって決まる.さらに,
  \[
    \frac{\theta}{2\pi}
  \]
  が有理数であるときには,ある正の整数$T$と整数$k$について
  \[
    T\theta=2k\pi
  \]
  とすることができる.このとき,
  \[
    \cos(n+T)\theta=\cos n\theta
  \]
  および
  \[
    \sin(n+T)\theta=\sin n\theta
  \]
  が成り立つため,
  \[
    a_{n+T}=a_n
  \]
  となる.したがって,この場合には数列そのものが周期的になる.

  一方,
  \[
    \rho\neq1
  \]
  の場合には,\,$\cos n\theta$や$\sin n\theta$が周期的に振動していても,\,$\rho^n$によって振幅が変化するため,一般には数列そのものは周期的にはならない.すなわち,虚数解をもつ漸化式では,\,$\rho$が数列の振幅の変化を,\,$\theta$がその振動の周期的な性質を決めていると考えることができる.